% ============================================================
%  EXAMEN TEMA 5: ESTRUCTURA ATÓMICA
%  Curso 4to Año — Química
%  Duración: 90 minutos
%  Temas: Modelos atómicos, partículas subatómicas, configuración
%         electrónica y números cuánticos
%  Autor: Ing. Gabriel Astudillo
%  Nota: enunciado limpio (sin pistas ni desarrollo). Las
%  soluciones están en la "Clave de Respuestas" al final.
% ============================================================

\documentclass[12pt, letterpaper]{article}

% ── Paquetes ─────────────────────────────────────────────────

\usepackage{mathwizards-guia}
\mwguia{Examen — Tema 5: Estructura Atómica}{Ing. Gabriel Astudillo $\cdot$ 4to Año}

\begin{document}
% ============================================================

% ── Portada / Título ─────────────────────────────────────────
\mwportada{Examen del Tema 5}{Estructura Atómica}{Modelos Atómicos, Partículas Subatómicas, Configuración Electrónica y Números Cuánticos}

\vspace{4pt}
\begin{cuadroinfo}
  \small
  \textbf{Nombre:} \rule{0.55\linewidth}{0.4pt} \quad \textbf{Fecha:} \rule{0.15\linewidth}{0.4pt} \\[6pt]
  \textbf{Tiempo disponible:} 90 minutos \quad$\bullet$\quad \textbf{Puntaje total:} 100 puntos \\[4pt]
  \textbf{Instrucciones:}
  \begin{enumerate}[leftmargin=*]
    \item Lee cada ejercicio con calma antes de resolverlo.
    \item Muestra \emph{todos} tus pasos: configuración electrónica antes de responder.
    \item Recuerda: $A - Z$ da los neutrones y $Z$ identifica al elemento.
    \item Aplica el principio de Aufbau y la regla de Hund.
    \item No se permite el uso de calculadora.
    \item Escribe con letra clara y revisa tu examen antes de entregarlo.
  \end{enumerate}
\end{cuadroinfo}

\vspace{8pt}

% ============================================================
\section{Modelos Atómicos y Partículas Subatómicas (34 puntos)}
% ============================================================

\begin{enumerate}[leftmargin=*]
  \item \textbf{(10 pts)} Un átomo de fósforo tiene $Z = 15$ y $A = 31$.
        \begin{enumerate}[label=\alph*)]
          \item \textquestiondown Cuántos protones, neutrones y electrones tiene el átomo neutro?
          \item Escribe su notación ${}^{A}_{Z}X$.
        \end{enumerate}

  \item \textbf{(8 pts)} Compara los isótopos ${}^{35}_{17}\text{Cl}$ y ${}^{37}_{17}\text{Cl}$. \textquestiondown Qué tienen en común y en qué se diferencian? Calcula los neutrones de cada uno.

  \item \textbf{(8 pts)} Indica cuántos protones y electrones tienen:
        \[
          a)\ \text{Mg}^{2+} \ (Z = 12) \qquad b)\ \text{O}^{2-} \ (Z = 8)
        \]
        Además, indica si cada uno es un \textbf{catión} o un \textbf{anión}.

  \item \textbf{(8 pts)}
        \begin{enumerate}[label=\alph*)]
          \item Ordena cronológicamente los modelos atómicos de Dalton, Thomson, Rutherford y Bohr.
          \item \textquestiondown Qué descubrió Rutherford con su experimento de la lámina de oro?
        \end{enumerate}
\end{enumerate}

\newpage

% ============================================================
\section{Configuración Electrónica (34 puntos)}
% ============================================================

\begin{enumerate}[leftmargin=*]
  \item \textbf{(8 pts)} Escribe la configuración electrónica de:
        \[
          a)\ \text{Cl} \ (Z = 17) \qquad b)\ \text{Ca} \ (Z = 20)
        \]

  \item \textbf{(8 pts)} Escribe la configuración electrónica y determina el número de electrones de valencia de:
        \[
          a)\ \text{azufre} \ (Z = 16) \qquad b)\ \text{argón} \ (Z = 18)
        \]

  \item \textbf{(9 pts)} Escribe la configuración electrónica del fósforo ($Z = 15$) y representa los orbitales $3p$ usando cajas y flechas, aplicando la \textbf{regla de Hund}.

  \item \textbf{(9 pts)} El ion $\text{S}^{2-}$ proviene del azufre ($Z = 16$). Escribe su configuración electrónica y compárala con la del argón ($Z = 18$). \textquestiondown Qué observas?
\end{enumerate}

\newpage

% ============================================================
\section{Números Cuánticos y Tabla Periódica (32 puntos)}
% ============================================================

\begin{enumerate}[leftmargin=*]
  \item \textbf{(8 pts)} Determina los cuatro números cuánticos $(n, l, m_l, m_s)$ del último electrón del fósforo ($Z = 15$).

  \item \textbf{(8 pts)} El calcio tiene $Z = 20$. Escribe su configuración electrónica y determina en qué \textbf{periodo} y \textbf{grupo} de la tabla periódica se encuentra.

  \item \textbf{(8 pts)} Un elemento está en el \textbf{periodo 3, grupo 16}. \textquestiondown Cuál es su configuración electrónica y su número atómico $Z$? \textquestiondown De qué elemento se trata?

  \item \textbf{(8 pts)} El oxígeno tiene $Z = 8$.
        \begin{enumerate}[label=\alph*)]
          \item Escribe su configuración electrónica.
          \item Aplicando la regla de Hund, \textquestiondown cuántos electrones desapareados tiene en el subnivel $2p$?
        \end{enumerate}
\end{enumerate}

\newpage

% ============================================================
\ifmwclaves
  \section{Clave de Respuestas}
  % ============================================================

  \subsection*{Sección 1: Modelos Atómicos y Partículas Subatómicas}

  \begin{enumerate}[leftmargin=*, label=\textbf{\arabic*.}]
    \item a) Protones $= Z = 15$; neutrones $= A - Z = 31 - 15 = 16$; electrones $= 15$ (átomo neutro).
          \quad b) ${}^{31}_{15}\text{P}$.

    \item Ambos tienen el mismo $Z = 17$: 17 protones y 17 electrones. Difieren en los neutrones:
          ${}^{35}\text{Cl}$ tiene $35 - 17 = 18$ neutrones y ${}^{37}\text{Cl}$ tiene $37 - 17 = 20$ neutrones.

    \item a) $\text{Mg}^{2+}$: 12 protones y $12 - 2 = 10$ electrones (perdió 2); es un \textbf{catión}.
          \quad b) $\text{O}^{2-}$: 8 protones y $8 + 2 = 10$ electrones (ganó 2); es un \textbf{anión}.

    \item a) Dalton (1803) $\rightarrow$ Thomson (1897) $\rightarrow$ Rutherford (1911) $\rightarrow$ Bohr (1913).
          \quad b) Rutherford descubrió el \textbf{núcleo} atómico: pequeño, denso y positivo. En su experimento, la mayoría de las partículas alfa atravesaban la lámina de oro sin desviarse (el átomo es mayormente vacío), pero algunas rebotaban al chocar con el núcleo.
  \end{enumerate}

  \newpage

  \subsection*{Sección 2: Configuración Electrónica}

  \begin{enumerate}[leftmargin=*, label=\textbf{\arabic*.}]
    \item a) Cl ($17$ e$^-$): $1s^{2}\,2s^{2}\,2p^{6}\,3s^{2}\,3p^{5}$.
          \quad b) Ca ($20$ e$^-$): $1s^{2}\,2s^{2}\,2p^{6}\,3s^{2}\,3p^{6}\,4s^{2}$.

    \item a) Azufre ($16$ e$^-$): $1s^{2}\,2s^{2}\,2p^{6}\,3s^{2}\,3p^{4}$ $\Rightarrow$ valencia $= 2 + 4 = 6$.
          \quad b) Argón ($18$ e$^-$): $1s^{2}\,2s^{2}\,2p^{6}\,3s^{2}\,3p^{6}$ $\Rightarrow$ valencia $= 2 + 6 = 8$.

    \item Fósforo ($15$ e$^-$): $1s^{2}\,2s^{2}\,2p^{6}\,3s^{2}\,3p^{3}$.
          \\[2pt]
          Los tres electrones $3p$ se colocan \emph{sueltos} (regla de Hund), con espines paralelos:
          \[
            3p:\ \boxed{\uparrow}\ \boxed{\uparrow}\ \boxed{\uparrow}
          \]

    \item El azufre neutro tiene 16 electrones; el ion $\text{S}^{2-}$ ganó 2 electrones, así que tiene 18:
          \[
            \text{S}^{2-}:\ 1s^{2}\,2s^{2}\,2p^{6}\,3s^{2}\,3p^{6}.
          \]
          Esta es \textbf{exactamente} la configuración del argón ($Z = 18$): se dice que son \textbf{isoelectrónicos}.
  \end{enumerate}

  \newpage

  \subsection*{Sección 3: Números Cuánticos y Tabla Periódica}

  \begin{enumerate}[leftmargin=*, label=\textbf{\arabic*.}]
    \item Configuración del fósforo: $1s^{2}\,2s^{2}\,2p^{6}\,3s^{2}\,3p^{3}$. El último subnivel es $3p$, con $n = 3$ y $l = 1$.
          Los tres electrones $3p$ van sueltos en los orbitales $-1, 0, +1$ con espín $+\tfrac{1}{2}$.
          Último electrón: $(n, l, m_l, m_s) = (3, 1, +1, +\tfrac{1}{2})$.

    \item Configuración: $1s^{2}\,2s^{2}\,2p^{6}\,3s^{2}\,3p^{6}\,4s^{2}$.
          Último nivel: $4$ $\Rightarrow$ \textbf{periodo 4}. Electrones de valencia: $4s^{2}$ $\Rightarrow$ 2 electrones $\Rightarrow$ \textbf{grupo 2} (alcalinotérreos).

    \item Periodo 3 y grupo 16 $\Rightarrow$ 6 electrones de valencia en el nivel 3: $3s^{2}\,3p^{4}$.
          Configuración completa: $1s^{2}\,2s^{2}\,2p^{6}\,3s^{2}\,3p^{4}$, con $Z = 16$. Se trata del \textbf{azufre (S)}.

    \item a) Oxígeno ($8$ e$^-$): $1s^{2}\,2s^{2}\,2p^{4}$.
          \quad b) En $2p^{4}$, primero se llenan los tres orbitales con un electrón cada uno ($-1, 0, +1$) y el cuarto se aparea:
          \[
            2p:\ \boxed{\uparrow\downarrow}\ \boxed{\uparrow}\ \boxed{\uparrow}
          \]
          Por lo tanto, el oxígeno tiene \textbf{2 electrones desapareados} en $2p$.
  \end{enumerate}

\fi
\end{document}
